% LSPV-3 REPORT::VALSUBCUBEFW1::2026-10-14T00:38:20+09:00
% !TEX program = xelatex
\documentclass[11pt]{article}
\usepackage[T1]{fontenc}
\usepackage{lmodern,amsmath,amssymb,amsthm,booktabs,tabularx,xurl,hyperref}
\usepackage[margin=0.86in]{geometry}
\newcommand{\EXACT}{\textbf{PROVED/EXACT}}
\newcommand{\KNOWN}{\textbf{KNOWN THEOREM}}
\newcommand{\COMP}{\textbf{COMPUTATIONAL}}
\newcommand{\COND}{\textbf{CONDITIONAL}}
\newcommand{\CAND}{\textbf{CANDIDATE}}
\newcommand{\OPEN}{\textbf{OPEN}}
\newcommand{\RET}{\textbf{FALSIFIED/RETIRED}}
\newtheorem{theorem}{Theorem}
\hypersetup{colorlinks=true,urlcolor=blue,linkcolor=blue}
\setlength{\emergencystretch}{3em}
\title{VALSUBCUBEFW1: Fixed-Valuation Subfield Cube Selector\\
Red-Team Report}
\author{UNIVERSAL\_CORE\_REDTEAM / LSPV-3}
\date{14 October 2026}

\begin{document}
\maketitle

\begin{center}
\fbox{\parbox{0.95\linewidth}{
\textbf{Executive status: P versus NP remains OPEN.}
CUBEROOTSEL1 is correct and is promoted here to \EXACT{} after independent
red-team review.  More strongly, the apparent $6\mid d$ escape is not a
sufficient repair.  For every fixed $a=\nu_3(d)$, an explicit subfield of
$\mathbb F_{2^d}$ deterministically supplies a cubic nonresidue and an
Adleman--Manders--Miller-style target-only selector in polynomial time.  In
particular the genuine square-cotree case $K_{7,7}$ has $d=36$ and is easy by
the embedded $\mathbb F_{2^{18}}$.  This is an upper-bound firewall, not an
unrestricted lower bound.}}
\end{center}

\section{Claim card}
\begin{center}
\begin{tabularx}{0.98\linewidth}{p{0.25\linewidth}X}
\toprule
Claim ID & VALSUBCUBEFW1 \\
Status before run & $6\nmid d$ selector \CAND; $6\mid d$ deterministic cube-root cost \OPEN \\
Target & Falsify divisibility by six as a sufficient target-only-selector repair \\
Quantifiers & For each fixed $a\ge1$, every even $d$ with $\nu_3(d)=a$, every explicit irreducible degree-$d$ modulus, and every target $b\in\mathbb F_{2^d}$ \\
Machine/model & Uniform deterministic random-access computation with polynomial workspace; $a$ is fixed for the family \\
Resource & Parsing, subfield discovery, nonresidue search, exponentiation, correction search, verification, and Tseitin reconstruction \\
SAME\_REALIZATION & The same typed field, target $b$, selected root, charge vector, $K_{s,s}$ cotree, and reconstructed edge assignment \\
Falsifier & One infinite $6\mid d$ typed subtype with a source-independent polynomial selector \\
\bottomrule
\end{tabularx}
\end{center}

\section{Primary literature and transfer classification}
Adleman, Manders, and Miller introduced finite-field root extraction with a
supplied nonresidue \cite{AMM}; the detailed cubic correction structure was
later made explicit by Cao, Sha, and Fan \cite{Cao}.  Sze gives deterministic
root extraction without a supplied nonresidue under additional arithmetic
hypotheses and explicitly warns that efficiency is conditional on the
factor/roots-of-unity regime \cite{Sze}.  Barreto--Voloch exploit exponent
periodicity in complementary extension regimes \cite{BV}.

The classification is exact.  Finite-field root extraction with a derived
nonresidue is a \textbf{directly transferable theorem mechanism}: its field,
input and operations are exactly those of the residual.  Cryptographic
root-extraction hardness is a \textbf{misleading analogy / failed transfer}:
no hardness assumption is admissible here, and the typed subfield is instead
an algorithmic weakness.  General deterministic root finding without the
stated arithmetic hypotheses remains \OPEN{} and is not inferred from these
papers.

\section{Independent review of CUBEROOTSEL1}
Let $d\ge4$ be even with $6\nmid d$, put
$q=2^d$ and $m=(q-1)/3$.  Then $3\nmid m$.  For
$e=3^{-1}\pmod m$, every nonzero cube $b$ lies in the subgroup of order $m$,
so
\[
  (b^e)^3=b^{3e}=b^{1+km}=b.
\]
Zero is handled separately, and recomputation rejects off-image targets.
The proof, bit-length accounting, affine-conjugation rule, and executable
audit in CUBEROOTSEL1 are correct.  CUBEROOTSEL1 is therefore \EXACT{} after
this independent review; its finite measurements remain \COMP{}.

\section{Fixed-valuation subfield selector theorem}
\begin{theorem}[fixed-$3$-adic-valuation selector]
Fix $a\ge1$.  There is a uniform deterministic algorithm $R_a$ with runtime
polynomial in $d$ such that, for every even $d$ satisfying $\nu_3(d)=a$, every
explicit irreducible $f\in\mathbb F_2[X]$ of degree $d$, and every
$b\in K=\mathbb F_2[X]/(f)$, the algorithm returns an $x$ with $x^3=b$ iff
$b$ is a cube, and otherwise rejects.  The input contains no source preimage
and no nonresidue advice.
\end{theorem}

\begin{proof}
Set
\[
 r=2\cdot3^a,\qquad q=2^d,
 \qquad q-1=3^{a+1}m,
 \qquad 3\nmid m.
\]
Because $d/3^a$ is even, $r\mid d$.  The lifting-the-exponent identity gives
\[
 \nu_3(2^d-1)=1+\nu_3(d)=a+1
 =\nu_3(2^r-1).
\]
Thus $K$ contains the unique subfield $K_0\cong\mathbb F_{2^r}$ and $K_0^*$
has the full $3$-primary order of $K^*$.

From the polynomial-basis representation of $K$, construct the matrix of the
$\mathbb F_2$-linear map $z\mapsto z^{2^r}+z$.  Gaussian elimination returns
an $r$-element basis of its kernel, namely $K_0$.  Since $a$ is fixed,
enumerating the $2^r$ elements is a family-dependent constant.  Test each
nonzero $c$ by
\[
 c^{(q-1)/3}\ne1.
\]
Such a $c$ exists: a generator of $K_0^*$ is not a cube in $K$ because the
cofactor $(q-1)/(2^r-1)$ is prime to $3$.

For nonzero promised cube $b$, compute $e=3^{-1}\pmod m$, $y=b^e$, and
$\eta=c^m$.  The element $\eta$ has order $3^{a+1}$.  Because $b$ is a cube,
$y^3/b$ lies in $\langle\eta^3\rangle$.  Therefore exactly one of the
$3^a$ candidates
\[
 x_j=y\eta^j,\qquad 0\le j<3^a,
\]
satisfies $x_j^3=b$.  Test the candidates and return the successful one;
return zero on $b=0$ and reject if none succeeds.  Subfield linear algebra,
enumeration for fixed $a$, repeated squaring, candidate testing, and output
all cost $\operatorname{poly}(d)$ bit operations under the declared field
representation.  No source witness is accessed.
\end{proof}

The theorem is \CAND{} pending a second independent review, as required for
project promotion.  Its proof is elementary and does not assume GRH,
randomness, ETH, cryptographic hardness, or a root oracle.

\section{Square-cotree corollary}
For $K_{s,s}$ the cotree dimension is $d=(s-1)^2$.  Fix
$h=\nu_3(s-1)\ge1$.  Then $\nu_3(d)=2h$ is fixed and the theorem applies.
In particular, $s=6u+1$ with $3\nmid u$ gives an infinite family with
$d=(6u)^2$, and its selector derives a nonresidue from the fixed subfield
$\mathbb F_{2^{18}}$.  The first member $s=7$ has $d=36$ and unbounded
extensions retain the same target-only algorithm.  Hence the condition
$6\mid d$, even together with square cotree geometry and unbounded
treewidth, is \RET{} as a sufficient hardness repair.

This does not cover a single family in which $\nu_3(d)$ grows without bound.
It also does not show that such a family is hard.  It merely sharpens the
next hostile control: any proposed cube residual must defeat typed subfield
discovery and nonresidue-assisted correction on its actual degrees.

\section{Bounded reproducible audit}
The dependency-free program
\path{experiments/review/fixed_valuation_cube_selector_redteam.py} ran in
1.702 seconds with seed 38142026.  It used exact polynomial-basis arithmetic
and passed 7,659 assertions with zero failures:
\begin{itemize}
\item degrees $d=6,12,18,36$ with $\nu_3(d)=1,1,2,2$;
\item derived subfields of degrees $6,6,18,18$ and no supplied nonresidue;
\item 5,184 promised-image targets and 1,368 off-image rejections;
\item 3,485 selected roots differing from the erased source root;
\item 256 public scalar-affine conjugation controls;
\item 192 complete $K_{7,7}$ Tseitin witness reconstructions at $d=36$;
\item even/odd charge, malformed subfield degree, irreducibility, target
verification, and off-image controls.
\end{itemize}
The finite audit is \COMP{}.  It checks only $a=1,2$; the general fixed-$a$
claim rests on the symbolic proof, not runtime extrapolation.  Raw output is
\path{results/raw/review/2026-10-14_0038_fixed_valuation_cube_selector.json}.
The code SHA-256 is
\texttt{fe6a33517c02d112b0f0acd892d659d81a996578ef20686c63e90f7dcfc7c7d1}.

\section{Hostile controls and barrier audit}
The selector has full target access, polynomial workspace, explicit field
arithmetic, and charged subfield discovery, nonresidue search, correction,
verification, forest reconstruction, and output.  It is source-independent,
survives public affine phase/scaling changes, and rejects off-image targets.
Horn, Krom, affine Gaussian elimination, bounded-treewidth, products,
renaming, complementation, and padding remain easy-family controls from the
verified predecessor suite; none is used to infer this theorem.

The firewall relativizes and algebrizes: semantic or algebraic oracle access
can only make the displayed selector easier.  Natural Proofs is not invoked.
Restricted DNNF, PC/PCR, Resolution, pebbling, one-tape, circuit, or
communication lower bounds supply no same-realization runtime lower bound.
DEAD/COPY/DELAY wrappers do not remove the selector; INTERP/RECOMP can execute
it directly.  Program length is distinct from workspace, and no fixed
percentage overhead is promoted to an exponent gap.

For the post-quantum ledger, the finite-field multiplication and subfield
object are an \textbf{EXACT OBJECT/THEOREM TRANSFER}.  Schur-square,
quadratic-relation, Hilbert/syzygy, collision-entropy, and LWE-style external
extraction ideas are \textbf{FAILED TRANSFERS} to unrestricted runtime here:
the algebraic structure is explicitly exploitable by the selector.  Generic
binding language is only an \textbf{ANALOGY/CANDIDATE} and supplies no
unconditional lower bound.

\section{CORPUS REVIEWED}
The run surveyed the current GitHub tree; root \path{main.tex};
\path{AGENTS.md}, \path{CODEX.md}, the complete PNP research skill and all
seven task specifications; newest MAIN PROMISEGATE1 and historical
CUBEPERMFW1; AFFINE-RESIDUAL-REDIRECT1, EARLY22-INTEGRATION1 and
SINGLE22-SEPARATION1; AFFINE-RESIDUAL-RANK1; the CCSPC, HEC, G-PPC-HEC and
CRS-PPC causal components; current PLANTEDCUBEFW1, CUBEROOTSEL1,
CHECKABLEPERM1 and relevant LSPV-3 firewall reports; and the exact predecessor
code/raw output.  The immediately preceding LSPV-2 report has a verified
header, substantive body and final sentinel, so no write recovery was needed.

\section{DUPLICATION / CONTRADICTION AUDIT}
Repository searches for fixed $3$-adic valuation, subfield cubic nonresidue,
AMM cube correction, and unbounded valuation found no prior project theorem
matching VALSUBCUBEFW1.  CUBEROOTSEL1 handles only $\nu_3(d)=0$ by exponent
inversion.  PLANTEDCUBEFW1 uses a source section.  CUBEPERMFW1 handles odd
$d$.  The present target-only subfield selector is therefore new to the
corpus and contradicts no valid prior theorem.

\section{WHAT CHANGED RELATIVE TO ROOT MAIN}
\begin{itemize}
\item CUBEROOTSEL1 is independently verified and may be promoted from
\CAND{} to \EXACT{} by synthesis.
\item The root statement that $6\mid d$ is an open algorithmic case remains
literally true in full generality, but is too coarse: every fixed
$\nu_3(d)$ stratum has a polynomial selector.
\item The repair ``require $6\mid d$'' is retired as sufficient, including
an infinite genuine $K_{s,s}$ square-cotree subtype.
\item The unrestricted lower-bound arrow and P versus NP status are unchanged.
\end{itemize}

\section{ADOPTED / REJECTED / RETIRED PRIOR COMPONENTS}
\textbf{Adopted:} strict typed-field recognition, polynomial-basis
arithmetic, target verification, the exact forest lift, and all valid
exact-$(2,2)$ compiler/proof-system results at their restricted scope.
\textbf{Rejected:} an unstated nonresidue, probabilistic expected time, or a
cryptographic root-extraction assumption as an unconditional selector.
\textbf{Retired:} divisibility by six, square cotree geometry, and erased
source roots as a sufficient joint obstruction on fixed-valuation strata.

\section{STALE OR OVERSTATED CLAIMS FOUND}
The phrase ``remaining $6\mid d$ cube subtype'' in CUBEROOTSEL1 is not false,
but it overstates how much of that region remains algorithmically unresolved.
Only strata with unbounded $\nu_3(d)$ evade this fixed-subfield enumeration;
even there, no hardness conclusion follows.  No valid compiler, affine
elimination, projected-rank, or restricted proof-system theorem is affected.

\section{Status and handoff}
\begin{center}
\begin{tabularx}{0.98\linewidth}{p{0.39\linewidth}p{0.20\linewidth}X}
\toprule
Claim & Status & Scope \\
\midrule
CUBEROOTSEL1 image exponent theorem & \EXACT & Independent red-team proof review complete \\
Fixed-$\nu_3(d)$ subfield selector & \CAND & Complete proof here; second review pending \\
Finite field/affine/$K_{7,7}$ audit & \COMP & 7,659 assertions, zero failures \\
$6\mid d$ as sufficient cube repair & \RET & Infinite fixed-valuation square-cotree subtype is easy \\
Unbounded-$\nu_3(d)$ target selector & \OPEN & Fixed-subfield enumeration no longer polynomial uniformly \\
Unrestricted all-program lower bound & \OPEN & No model-matched transfer \\
P versus NP separation & \OPEN & Load-bearing implication remains open \\
\bottomrule
\end{tabularx}
\end{center}

\textbf{Single next target.}  Audit the unbounded-$\nu_3(d)$ regime using a
uniform tower algorithm: derive the required $3$-primary generator without
enumerating $2^{2\cdot3^{\nu_3(d)}}$ elements, and test Sze-style
nonresidue-free extraction on the exact typed representation.  Failure of
known algorithms is not a lower bound.

\bigskip
\noindent P versus NP remains open.

\begin{thebibliography}{9}
\bibitem{AMM}
L. Adleman, K. Manders, and G. Miller,
\emph{On Taking Roots in Finite Fields}, 18th FOCS, 1977, 175--178.
\bibitem{Cao}
Z. Cao, Q. Sha, and X. Fan,
\emph{Adleman--Manders--Miller Root Extraction Method Revisited},
arXiv:1111.4877, 2011.
\bibitem{Sze}
T.-W. Sze,
\emph{On Taking $r$-th Roots without $r$-th Nonresidues over Finite Fields
and Its Applications}, arXiv:1105.5852, 2011.
\bibitem{BV}
P. S. L. M. Barreto and J. F. Voloch,
\emph{Efficient Computation of Roots in Finite Fields},
Designs, Codes and Cryptography 39 (2006), 275--280.
\end{thebibliography}

\end{document}
% END_OF_VERIFIED_REPORT::LSPV-3::VALSUBCUBEFW1
